10.8 Sum and Difference Formulas

1. The Formulas
$$ \sin(a \pm b) = \sin a \cos b \pm \cos a \sin b $$
$$ \cos(a \pm b) = \cos a \cos b \mp \sin a \sin b $$
$$ \tan(a \pm b) = \frac{\tan a \pm \tan b}{1 \mp \tan a \tan b} $$

2. Exact Values (Ex. 1)
Break angle into two known angles, then apply formula.
$$ \sin 15° = \sin(60°-45°) = \frac{\sqrt{6}-\sqrt{2}}{4} $$
$$ \tan\frac{7\pi}{12} = \tan\!\left(\frac{\pi}{3}+\frac{\pi}{4}\right) = \frac{\sqrt{3}+1}{1-\sqrt{3}} = -2-\sqrt{3} $$

3. Given sin/cos, Find cos(a-b) (Ex. 2)
Use Pythagorean identity to find missing sides first.
$$ \cos a = -\frac{4}{5} \Rightarrow \sin a = -\frac{3}{5} \text{ (QIII)} $$
$$ \sin b = \frac{5}{13} \Rightarrow \cos b = \frac{12}{13} \text{ (QI)} $$
$$ \cos(a-b) = \cos a\cos b + \sin a\sin b = -\frac{48}{65}+\frac{-15}{65} = -\frac{63}{65} $$

4. Simplifying (Ex. 3)
$$ \cos(x+\pi) = \cos x\cos\pi - \sin x\sin\pi = -\cos x $$

5. Solving Equations (Ex. 4)
$$ \sin\!\left(x+\frac{\pi}{3}\right)+\sin\!\left(x-\frac{\pi}{3}\right) = 1 $$
Expand both, terms cancel: $\sin x = 1$
$$ x = \frac{\pi}{2} $$
